% !TEX program = pdflatex
\documentclass[11pt,a4paper]{article}

\usepackage[utf8]{inputenc}
\usepackage[T1]{fontenc}
\usepackage{lmodern}
\usepackage[margin=2.4cm]{geometry}
\usepackage{amsmath,amssymb,amsthm}
\usepackage{booktabs}
\usepackage{enumitem}
\usepackage{microtype}
\usepackage{hyperref}
\usepackage{xcolor}

\hypersetup{
  colorlinks=true,
  linkcolor=black,
  citecolor=black,
  urlcolor=blue!50!black
}

\newtheorem{theorem}{Theorem}
\newtheorem{proposition}[theorem]{Proposition}
\theoremstyle{definition}
\newtheorem{definition}[theorem]{Definition}
\newtheorem{hypothesis}[theorem]{Hypothesis}
\theoremstyle{remark}
\newtheorem{remark}[theorem]{Remark}

\newcommand{\Lean}{\textsc{Lean}\,4}
\newcommand{\CA}{\ensuremath{[\mathrm{C}\!\to\!\mathrm{A}]}}
\newcommand{\tagA}{\ensuremath{[\mathrm{A}]}}
\newcommand{\tagB}{\ensuremath{[\mathrm{B}]}}
\newcommand{\tagC}{\ensuremath{[\mathrm{C}]}}
\newcommand{\FeedIn}{\mathtt{core6PathFeedInSet}}
\newcommand{\Residual}{\mathtt{core6PathResidualSet}}
\newcommand{\Umap}{U}

\title{%
  From the Collatz Suspicion to a Residual Reduction:\\
  A Letter on the Pre119--Core6 Formal Path
}
\author{%
  Thomas Hofbauer\\
  \small Bamberg / Akademie Olympia\\
  \small\texttt{Kepler--Hurwitz} formalization programme
}
\date{27 July 2026}

\begin{document}
\maketitle

\begin{abstract}
This note is addressed, with respect, to \textbf{Lothar Collatz},
to \textbf{Riho Terras}, and to the many mathematicians who shaped the
modern language of the $3x+1$ problem.
It recounts---without claiming a proof---how a contemporary
machine-checked investigation isolates the classical suspicion inside a
precise residual set.
On the Core6 cylinder architecture we obtain a complete algebraic
classification of the \emph{Core6-internal} first-hit feed-in channel and a
formal equivalence
\[
  \mathtt{ResidualFeedInGoal}
  \;\Longleftrightarrow\;
  \mathtt{ReachabilityFeedInGoal}.
\]
Thus the global dynamical difficulty is localized, not dissolved.
The epistemic status remains deliberately open: \texttt{ClaimsFreeze}
is false; Collatz's conjecture is unproved.
\end{abstract}

\tableofcontents
\vspace{1em}

%==============================================================================
\section{An address to those who came before}
%==============================================================================

Dear Professor Collatz,
dear Dr.~Terras,
and dear colleagues along the long $3x+1$ road---
Everett, Kakutani, Ulam, Lagarias, Conway, Wirsching, Korec, Crandall,
Kontorovich, Tao, and so many unnamed computers and computers of computers---

We write not to announce a solution, but to report a \emph{localization}.
Your suspicion---that every positive integer, under the map that doubles down
on odds and halves on evens, eventually reaches the cycle
$4\to 2\to 1$---remains open.
What we can offer, after a day of concentrated formal work inside the
\texttt{Kepler--Hurwitz} Lean library, is a sharper statement of
\emph{where} the remaining difficulty lives.

The classical problem asks for a global dynamical fact.
Our contribution is architectural:
we carve the odd Syracuse dynamics along a fixed valuation word
(\texttt{Core6}), prove what can be proved inside that word, refute what
cannot, and reduce the rest to a residual set that is logically equivalent
to full expanding-mass reachability.

If this letter has one moral, it is the one Terras already taught the
subject with probabilistic clarity:
\begin{quote}
\emph{Almost all} behaviour is not \emph{all} behaviour;\\
and a beautiful average is not a uniform certificate.
\end{quote}
We try to keep that distinction machine-checkable.

%==============================================================================
\section{The suspicion, in classical language}
%==============================================================================

\subsection{Collatz's map}
Let
\[
  T(n) \;=\;
  \begin{cases}
    n/2, & n\text{ even},\\
    3n+1, & n\text{ odd}.
  \end{cases}
\]
The Collatz conjecture asserts that for every $n\in\mathbb{N}_{>0}$ there
exists $t\ge 0$ with $T^{\circ t}(n)=1$.
An equivalent and often cleaner formulation works with the
\emph{odd Syracuse map}
\[
  \Umap(n)
  \;=\;
  \frac{3n+1}{2^{v_2(3n+1)}},
\]
the first odd iterate after an odd kick.
Reachability of the trivial cycle is then equivalent to every odd start
eventually entering a contracting regime under iterates of $\Umap$.

\subsection{Voices on the path}
We name only a few of those whose work prepared the ground for any modern
attempt---including ours---to speak carefully about the problem.

\begin{description}[leftmargin=1.6cm,style=nextline]
  \item[Lothar Collatz (1910--1990).]
  The conjecture that bears his name: a child's game that became a
  century's embarrassment and invitation.
  We address him first because honesty toward the originator forbids
  packaging a localization as a proof.

  \item[Riho Terras.]
  Terras's density and stopping-time analysis remains the paradigmatic
  lesson that probabilistic almost-sure descent is not a uniform theorem.
  When we write ``Mr.~Taro'' in conversation, we mean Terras:
  the voice that taught the subject to distinguish heuristic measure from
  certificate.

  \item[C.~J.~Everett; Shizuo Kakutani; Stanis\l{}aw Ulam.]
  Early dynamical and probabilistic framings of the $3x+1$ iteration;
  the problem's migration from recreational oddity to ergodic curiosity.

  \item[Jeffrey C.~Lagarias.]
  The definitive surveys and the culture of precise partial results:
  congruence constraints, covering systems, computational frontiers,
  and the discipline of stating what is \emph{not} proved.

  \item[John H.~Conway.]
  Undecidability for generalized Collatz maps: a warning that formal
  caution is not pedantry.

  \item[G\"unter Wirsching; Ivan Korec; and the computational school.]
  Structure theorems, density estimates, and the vast experimental
  verification that makes every residual witness---such as our $31$---
  feel both familiar and sharp.

  \item[Terence Tao (2019).]
  Almost all orbits attain almost bounded values (logarithmic density).
  A magnificent external milestone, and still not the Collatz conjecture.
  We cite it as context, never as a discharged step of our Lean chain
  \cite{Tao2019}.
\end{description}

None of these names is invoked as co-author of the present formal package.
They are the audience to whom we owe an accurate status report.

%==============================================================================
\section{Our programme: Pre119 and the Core6 word}
%==============================================================================

\subsection{Epistemic tags}
Inside the project we refuse a single boolean ``proved / unproved'' for every
sentence.
Claims carry tags:
\begin{center}
\begin{tabular}{@{}lll@{}}
\toprule
Tag & Meaning & Promotion rule \\
\midrule
\tagA & Lean-formal, accepted after review & repository of record \\
\tagB & reproducible census / numerics & never alone upgrades to \tagA \\
\tagC & open hypothesis / interface & may become \CA \\
\CA & constructive Lean package, not yet \tagA & needs merge + separate promotion \\
\bottomrule
\end{tabular}
\end{center}
Global \texttt{ClaimsFreeze} remains \textbf{false}: we do not claim Collatz.

\subsection{The Core6 fiber}
A \emph{fiber word} is a finite sequence of positive $2$-adic valuations
recording the successive halvings after odd Syracuse kicks.
The distinguished prefix
\[
  \mathtt{core6}
  \;=\;
  [1,1,1,1,2,2]
\]
is a length-$6$ valuation word shared by a large coherent family of starts.
Appending a positive tail exponent $e$ yields the length-$7$ fiber
\[
  \mathtt{fiberE}(e)
  \;=\;
  \mathtt{core6}\mathbin{+\kern-0.2em+}[e].
\]
For each $e\ge 1$ there is a unique \emph{canonical base} $b_e$ in a complete
residue system modulo $2^{e+9}$, and an infinite arithmetic progression---the
\emph{canonical cylinder}
\[
  C_e
  \;=\;
  \bigl\{\, b_e + k\cdot 2^{e+9}
  \;\big|\;
  k\in\mathbb{N}\,\bigr\}.
\]
Dynamically we split
\begin{align*}
  \mathtt{expandingMass}
  &\;=\; C_1\cup C_2\cup C_3,\\
  \mathtt{contractingMass}
  &\;=\; \textstyle\bigcup_{e\ge 4} C_e.
\end{align*}
Heuristically, expanding channels tend to grow under a full block, while
contracting channels tend to shrink.
The classical reachability goal, in this language, is:

\begin{hypothesis}[ReachabilityFeedInGoal \tagC]
Every $n\in\mathtt{expandingMass}$ admits $t\ge 1$ with
$\Umap^{\circ t}(n)\in\mathtt{contractingMass}$.
\end{hypothesis}

This is still open.
What changed today is its \emph{logical location}.

%==============================================================================
\section{The formal path of the last weeks}
%==============================================================================

We summarize the stack that led to today's residual reduction.
All packages below are presently at status \CA{} (CI-green candidates), not yet
promoted repository-\tagA{} after merge.

\subsection{PR\,\#15 --- canonical seed lifting}
Unique bases, realization of $\mathtt{fiberE}(e)$, and infinite arithmetic
lifting for every $e\ge 4$:
every cylinder point is an infinite AP of genuine Core6-tail realizations.
Algebraically, the odd affine coefficient along $\mathtt{fiberE}(e)$ is
independent of~$e$ (the constant $2347$ in the formalization).

\subsection{PR\,\#16 --- static cylinder partition}
Pairwise disjointness of the cylinders~$C_e$, the expanding/contracting phase
partition of the Core6 cylinder, and a \emph{relative} dyadic density statement:
the contracting share among Core6 residues tends to $1/8$.
Explicit non-claim: this is \textbf{not} natural density on~$\mathbb{N}$, and
not Collatz.

\subsection{PR\,\#17 --- one-block feed-in (D2b freeze)}
After one full $\mathtt{fiberE}(e_0)$ block, the starts in~$C_{e_0}$ that land
immediately in contracting mass form a set
$\mathtt{oneBlockFeedInSet}(e_0)$.
This set is a disjoint union of arithmetic progressions indexed by target
channels $f\ge 4$, via a Cancel-by-$2$ congruence bridge and a unique residue
class~$\kappa$.
A positive witness is $1{,}246{,}239\in\mathtt{oneBlockFeedInSet}(1)$;
the universal one-block claim is formally false via $31\mapsto 137$.

\subsection{PR\,\#18 --- multi-block paths and residual reduction (D3.3 freeze)}
This is the subject of today's investigations.
Mathematical freeze head:
\[
  \texttt{24bd5c83718da11c9ec2bfa8962de331be117369}.
\]
Packaging with non-semantic Lean comments only:
\[
  \texttt{adb2dd84ea87473dd4e6c49ace50de15cb676297}.
\]

%==============================================================================
\section{Today's theorems, stated carefully}
%==============================================================================

\subsection{Fixed-path arithmetic progressions}
For a fixed nonempty channel path $p$ with valid labels, the set of starts
that realize~$p$ is exactly an arithmetic progression $\mathrm{AP}_p$
(obtained by iterated dyadic lifting, not by an appeal to classical CRT as a
black box).

\subsection{First-hit packaging}
A \emph{contracting first-hit path} starts in an expanding channel, stays
expanding in every interior channel, and ends at the first contracting
channel.
Their union is
\[
  \FeedIn(e_0)
  \;=\;
  \bigcup_{\substack{p\text{ contracting}\\\text{first-hit from }e_0}}
  \mathrm{AP}_p.
\]
This is the Core6-\emph{internal} successful channel.

\subsection{Refutation of universal Core6-path coverage}
\begin{theorem}[$\neg\,\mathtt{Core6PathFeedInGoal}$ \CA]
It is not true that every start in $C_1\cup C_2\cup C_3$ realizes some
Core6-internal contracting first-hit path.
\end{theorem}
\begin{proof}[Proof sketch in the kernel]
Take $e_0=1$ and $n=31$.
Then $31\in C_1$, yet after one block $31\mapsto 137\notin\mathtt{core6Cylinder}$.
Any Core6-internal first-hit path would force the first block image back into
Core6; hence $31\notin\FeedIn(1)$.
\end{proof}
Thus the strategy
``every expanding start has a purely Core6-internal first-hit path''
is not merely unfinished---it is \textbf{false}.

\subsection{Cylinder dichotomy}
\begin{definition}
$\Residual(e_0) \;=\; C_{e_0}\setminus\FeedIn(e_0)$.
\end{definition}
\begin{theorem}[Dichotomy \CA]
$C_{e_0}=\FeedIn(e_0)\;\dot\cup\;\Residual(e_0)$,
and $\mathtt{oneBlockOffCore6Set}(e_0)\subseteq\Residual(e_0)$.
\end{theorem}
For $e_0=1$ the residual is nonempty ($31$) and the feed-in share is nonempty
($1{,}246{,}239$).

\subsection{The First-Hit bridge}
Membership in the feed-in set is not only arithmetic: it is dynamical.
\begin{theorem}[Feed-in bridge \CA]
If $n\in\FeedIn(e_0)$, then there exist $r\ge 1$ and
$t=7r$ with $\Umap^{\circ t}(n)\in\mathtt{contractingMass}$.
\end{theorem}
The Core6-internal share therefore needs \textbf{no further open dynamics}.

\subsection{The equivalence that localizes the difficulty}
\begin{theorem}[Residual reduction \CA]
\[
  \mathtt{ResidualFeedInGoal}
  \;\Longleftrightarrow\;
  \mathtt{ReachabilityFeedInGoal}.
\]
\end{theorem}
\begin{remark}
The reverse implication is immediate: residual starts lie in expanding mass.
The forward implication needs the dichotomy \emph{and} the First-Hit bridge.
The bare set equation $C=F\,\dot\cup\,R$ does not by itself move dynamics from
$F$ to contracting mass.
\end{remark}

In one sentence for Professor Collatz:
\begin{quote}
\emph{Everything that can be finished inside Core6-internal first-hit paths
is finished; everything that remains is exactly the residual, and the residual
is as hard as the whole expanding-mass reachability question.}
\end{quote}

%==============================================================================
\section{What the residual still hides}
%==============================================================================

The residual is not homogeneous.
At present we know only that immediate one-block Off-Core6 exits sit inside it.
Design architecture for a follow-up (not yet Lean, not a theorem) proposes
\[
  \Residual
  \;=\;
  \mathtt{FiniteBlockExit}
  \;\dot\cup\;
  \mathtt{ForeverExpandingBoundary},
\]
i.e.\ starts that leave Core6 after finitely many expanding blocks, versus
starts that remain forever in $C_1\cup C_2\cup C_3$ at every block boundary.
Re-entry after exit, and the exclusion or control of forever-expanding
boundary orbits, remain \tagC.

\begin{hypothesis}[ResidualFeedInGoal \tagC]
Every residual start eventually reaches contracting mass under odd Syracuse
iteration.
\end{hypothesis}
By the equivalence above, discharging this hypothesis would discharge
expanding-mass reachability in the Core6 language---still short of a full
Collatz theorem for all of~$\mathbb{N}$, but the natural next fortress.

%==============================================================================
\section{Governance: how we refuse to overclaim}
%==============================================================================

\begin{enumerate}[leftmargin=*]
  \item PR\,\#18 is under a \textbf{strict mathematical freeze} at
  \texttt{24bd5c8}. No new kernel mathematics belongs on that pull request.
  \item Status tags stay \CA{} until review, merge, and a \textbf{separate}
  promotion commit to repository-\tagA.
  \item D3.4 residual fine-structure and D4 residual dynamics belong on a
  follow-up branch
  (\texttt{cursor/core6-dynamic-d34-residual-\ldots}), not on the frozen tip.
  \item Natural density, global Collatz, and ``collapse'' rhetoric are
  explicit non-goals of the present package.
\end{enumerate}

This discipline is our debt to Conway's undecidability warning and to
Lagarias's survey culture: name the theorem you have, and name the theorem
you do not.

%==============================================================================
\section{A small diagram of the day}
%==============================================================================

\[
\begin{array}{c}
C_{e_0}\quad(1\le e_0\le 3)\\[0.4em]
\hline
\begin{array}{c|c}
\FeedIn(e_0)=\bigcup_p\mathrm{AP}_p
&
\Residual(e_0)\\[0.3em]
\text{First-Hit bridge: }t=7r
&
\text{open dynamical core}\\[0.2em]
\text{finished}
&
\equiv\ \mathtt{ReachabilityFeedInGoal}
\end{array}
\end{array}
\]

%==============================================================================
\section{Closing words}
%==============================================================================

Professor Collatz: your map is still unconquered.
What we can say, after formalizing the Core6 channel, is that one seductive
proof idea---universal coverage by Core6-internal first-hit paths---is dead,
and that the surviving difficulty has a clean name.

Dr.~Terras: your lesson stands.
We have converted a heuristic partition into Lean definitions and theorems,
but the residual still asks for dynamics that averages cannot supply.

To Lagarias, Tao, and the computational witnesses of decades:
thank you for the language in which an honest residual can be stated.

The conjecture remains a suspicion.
Our progress is that the suspicion now has a smaller, sharper address.

\vspace{1.2em}
\begin{flushright}
\textit{Bamberg, 27 July 2026}\\
Thomas Hofbauer\\
\texttt{Kepler--Hurwitz / Pre119--Core6}
\end{flushright}

%==============================================================================
\begin{thebibliography}{9}
\bibitem{Collatz}
L.~Collatz.
\newblock On the $3x+1$ problem (historical attribution; circulating notes and
correspondence, mid-20th century).
\newblock See surveys in~\cite{Lagarias2010}.

\bibitem{Terras1976}
R.~Terras.
\newblock A stopping time problem on the positive integers.
\newblock \emph{Acta Arith.} \textbf{30} (1976), 241--252.

\bibitem{Lagarias2010}
J.~C.~Lagarias (ed.).
\newblock \emph{The Ultimate Challenge: The $3x+1$ Problem}.
\newblock American Mathematical Society, 2010.

\bibitem{Conway1972}
J.~H.~Conway.
\newblock Unpredictable iterations.
\newblock In \emph{Proc.\ 1972 Number Theory Conf.}, Univ.\ of Colorado,
Boulder, 1972, 49--52.

\bibitem{Tao2019}
T.~Tao.
\newblock Almost all orbits of the Collatz map attain almost bounded values.
\newblock arXiv:1909.03562, 2019.

\bibitem{Wirsching1998}
G.~J.~Wirsching.
\newblock \emph{The Dynamical System Generated by the $3n+1$ Function}.
\newblock Lecture Notes in Mathematics \textbf{1681}, Springer, 1998.

\bibitem{KHPre119}
Akademie Olympia / Kepler--Hurwitz.
\newblock Pre119 Core6 formalization stack (PRs \#15--\#18),
math freeze heads
\texttt{3740183} (static),
\texttt{18e8747} (D2b),
\texttt{24bd5c8} (D3.3);
\texttt{ClaimsFreeze} false.
\end{thebibliography}

\end{document}
